\documentclass{beamer}

\usetheme{Madrid}

\usepackage{amsmath,amssymb,bm}
\usepackage{physics}
\usepackage{graphicx}
\usepackage{booktabs}

\title{Entropy-Based Effective Degrees of Freedom in Principal Component Analysis}
\author{Morten Hjorth-Jensen}
\date{\today}

\begin{document}

%------------------------------------------------
\begin{frame}
\titlepage
\end{frame}

%------------------------------------------------
\begin{frame}{Motivation}

In many applications of principal component analysis (PCA), one encounters situations where:

\begin{itemize}
    \item many eigenvalues are nonzero,
    \item but only a few contribute significantly,
    \item and one would like a quantitative measure of the
    \emph{effective dimensionality}.
\end{itemize}

\vspace{0.4cm}

Instead of asking:

\[
\text{``What is the exact rank?''}
\]

we ask:

\[
\text{``How many degrees of freedom are effectively active?''}
\]

\vspace{0.4cm}

This leads naturally to entropy-based measures of dimensionality.

\end{frame}

%------------------------------------------------
\begin{frame}{Principal Component Analysis}

Consider a covariance matrix

\[
C.
\]

Diagonalizing the covariance matrix gives

\[
C \mathbf{v}_i = \lambda_i \mathbf{v}_i,
\]

where:

\begin{itemize}
    \item $\lambda_i$ are the eigenvalues,
    \item $\mathbf{v}_i$ are the principal components.
\end{itemize}

\vspace{0.3cm}

The eigenvalues represent the variance carried by each principal direction.

\vspace{0.3cm}

The total variance is

\[
\sum_i \lambda_i.
\]

\end{frame}

%------------------------------------------------
\begin{frame}{Normalized Variance Fractions}

Define normalized variance fractions:

\[
p_i = \frac{\lambda_i}{\sum_j \lambda_j}
\]

\vspace{0.4cm}

These quantities satisfy

\[
\sum_i p_i = 1.
\]

\vspace{0.4cm}

Interpretation:

\begin{itemize}
    \item $p_i$ behaves like a probability distribution,
    \item each component measures the fraction of total variance,
    \item dominant PCA modes correspond to large $p_i$.
\end{itemize}

\end{frame}

%------------------------------------------------
\begin{frame}{Entropy of the Eigenvalue Spectrum}

The Shannon entropy associated with the PCA spectrum is

\[
S = -\sum_i p_i \log p_i
\]

\vspace{0.5cm}

This entropy measures:

\begin{itemize}
    \item how broadly distributed the variance is,
    \item how delocalized the spectrum is,
    \item how many modes contribute substantially.
\end{itemize}

\vspace{0.5cm}

Small entropy:
\begin{itemize}
    \item variance concentrated in a few modes.
\end{itemize}

Large entropy:
\begin{itemize}
    \item variance distributed among many modes.
\end{itemize}

\end{frame}

%------------------------------------------------
\begin{frame}{Effective Number of Degrees of Freedom}

The entropy itself is not directly an effective dimension.

\vspace{0.4cm}

Exponentiating the entropy gives:

\[
N_{\mathrm{eff}} = \exp\left(-\sum_i p_i \log p_i\right)
\]

\vspace{0.5cm}

This quantity is interpreted as:

\begin{block}{Interpretation}
The effective number of principal components contributing significantly to the data.
\end{block}

\vspace{0.4cm}

This is also called:

\begin{itemize}
    \item effective dimension,
    \item entropy dimension,
    \item effective rank.
\end{itemize}

\end{frame}

%------------------------------------------------
\begin{frame}{Limiting Case I: One Dominant Mode}

Suppose

\[
p_1=1,
\qquad
p_{i>1}=0.
\]

\vspace{0.4cm}

Then

\[
S=0.
\]

Hence

\[
N_{\mathrm{eff}}=e^0=1.
\]

\vspace{0.5cm}

Interpretation:

\begin{itemize}
    \item Only one effective degree of freedom exists.
    \item The data effectively lives in a one-dimensional manifold.
\end{itemize}

\end{frame}

%------------------------------------------------
\begin{frame}{Limiting Case II: Uniform Distribution}

Suppose $k$ modes contribute equally:

\[
p_i=\frac1k,
\qquad
i=1,\dots,k.
\]

Then

\[
S
=
-\sum_{i=1}^k
\frac1k
\log\frac1k.
\]

This simplifies to

\[
S=\log k.
\]

Therefore

\[
N_{\mathrm{eff}}
=
e^{\log k}
=
k.
\]

\vspace{0.4cm}

Thus the formula reproduces the intuitive number of active modes.

\end{frame}

%------------------------------------------------
\begin{frame}{Why Exponentiate the Entropy?}

Entropy itself is logarithmic.

\vspace{0.3cm}

Exponentiating converts the entropy into a quantity with units of:

\[
\text{``number of states''}.
\]

\vspace{0.4cm}

This is analogous to statistical mechanics:

\[
S = \log \Omega,
\]

where:

\begin{itemize}
    \item $S$ is entropy,
    \item $\Omega$ is the number of accessible microstates.
\end{itemize}

\vspace{0.4cm}

Thus

\[
e^S
\]

acts like an effective count of participating modes.

\end{frame}

%------------------------------------------------
\begin{frame}{Connection with Statistical Mechanics}

The PCA spectrum behaves similarly to occupation probabilities.

\vspace{0.4cm}

Analogy:

\begin{center}
\begin{tabular}{cc}
\toprule
Statistical Mechanics & PCA \\
\midrule
Microstates & Principal components \\
Occupation probabilities & Variance fractions $p_i$ \\
Entropy & Spectral entropy \\
Accessible states & Effective dimension \\
\bottomrule
\end{tabular}
\end{center}

\vspace{0.5cm}

This connection explains why entropy-based dimensionality measures are natural in physics.

\end{frame}

%------------------------------------------------
\begin{frame}{Participation Ratio}

Another commonly used measure is the participation ratio:

\[
N_{\mathrm{PR}} = \frac{1}{\sum_i p_i^2}
\]

\vspace{0.5cm}

Comparison:

\begin{itemize}
    \item entropy measure:
    \[
    e^{-\sum_i p_i\log p_i},
    \]
    
    \item participation ratio:
    \[
    \frac{1}{\sum_i p_i^2}.
    \]
\end{itemize}

\vspace{0.4cm}

Both quantify effective dimensionality.

\vspace{0.3cm}

The participation ratio penalizes dominant modes more strongly.

\end{frame}

%------------------------------------------------
\begin{frame}{Applications in Physics}

Entropy-based effective dimensions appear in many areas:

\vspace{0.3cm}

\begin{itemize}
    \item random matrix theory,
    \item localization theory,
    \item condensed matter physics,
    \item quantum entanglement spectra,
    \item many-body physics,
    \item neuroscience,
    \item machine learning,
    \item signal processing.
\end{itemize}

\vspace{0.4cm}

In many-body physics:

\begin{itemize}
    \item they measure the number of active collective modes,
    \item complexity of wave functions,
    \item degree of delocalization.
\end{itemize}

\end{frame}

%------------------------------------------------
\begin{frame}{Effective Rank of a Matrix}

In numerical linear algebra, the same quantity is often called the effective rank.

\vspace{0.4cm}

Given covariance matrix $C$:

\[
\mathrm{erank}(C)=\exp\left(-\sum_i p_i\log p_i\right)
\]

\vspace{0.5cm}

Advantages compared to ordinary rank:

\begin{itemize}
    \item robust against noise,
    \item insensitive to tiny eigenvalues,
    \item smoothly varying,
    \item physically meaningful.
\end{itemize}

\vspace{0.4cm}

Thus effective rank gives a continuous notion of dimensionality.

\end{frame}

%------------------------------------------------
\begin{frame}{Interpretation of Different Spectra}

\begin{center}
\begin{tabular}{cc}
\toprule
Spectrum & Effective Dimension \\
\midrule
One dominant eigenvalue & $\approx 1$ \\
Two equal dominant modes & $\approx 2$ \\
Broad spectrum & large \\
White noise spectrum & near full dimension \\
\bottomrule
\end{tabular}
\end{center}

\vspace{0.5cm}

The entropy-based measure compresses the entire spectrum into a single scalar quantity.

\end{frame}

%------------------------------------------------
\begin{frame}{Geometric Interpretation}

The effective dimension characterizes:

\begin{itemize}
    \item the intrinsic dimension of the data manifold,
    \item the number of active directions in feature space,
    \item the complexity of the covariance structure.
\end{itemize}

\vspace{0.4cm}

Low effective dimension:

\begin{itemize}
    \item strong correlations,
    \item compressible data,
    \item collective behavior.
\end{itemize}

\vspace{0.4cm}

High effective dimension:

\begin{itemize}
    \item weak correlations,
    \item distributed variance,
    \item noisy or high-complexity systems.
\end{itemize}

\end{frame}

%------------------------------------------------
\begin{frame}{Summary}

\begin{itemize}
    \item PCA eigenvalues define variance fractions
    \[
    p_i=\frac{\lambda_i}{\sum_j\lambda_j}.
    \]
    
    \item The Shannon entropy of the spectrum is
    \[
    S=-\sum_i p_i\log p_i.
    \]
    
    \item Exponentiating the entropy gives the effective number of degrees of freedom:
    \[
    N_{\mathrm{eff}}=e^S.
    \]
    
    \item This quantity measures the effective dimensionality of the data.
    
    \item Closely related concepts:
    \begin{itemize}
        \item participation ratio,
        \item effective rank,
        \item localization entropy,
        \item entanglement entropy.
    \end{itemize}
\end{itemize}

\end{frame}

%------------------------------------------------
\begin{frame}{References}

\footnotesize

\begin{thebibliography}{99}

\bibitem{Jolliffe}
I. T. Jolliffe,
\newblock {\it Principal Component Analysis},
Springer (2002).

\vspace{0.2cm}

\bibitem{RoyVetterli}
O. Roy and M. Vetterli,
\newblock ``The Effective Rank: A Measure of Effective Dimensionality,''
European Signal Processing Conference (2007).

\vspace{0.2cm}

\bibitem{Mezard}
M. Mezard, G. Parisi, and M. Virasoro,
\newblock {\it Spin Glass Theory and Beyond},
World Scientific (1987).

\vspace{0.2cm}

\bibitem{Anderson}
P. W. Anderson,
\newblock Phys. Rev. {\bf 109}, 1492 (1958).

\end{thebibliography}

\end{frame}

%------------------------------------------------
\end{document}
